\section*{Chapter \ref{ch:iv2:markets-and-microstructure} --- Markets and Microstructure}
\markboth{\protect\omTocMark{}{Answers: Markets and Microstructure}}{\protect\omTocMark{}{Answers: Markets and Microstructure}}

\begin{solution}{iq:iv2:markets-and-microstructure:1}
Your average price is about 100.0267, or about 2.67 cents above the mid-price of 100.00. The best offer becomes 100.03, for 200 shares, and the mid-price becomes 100.005.

You take all 100 at 100.02 and 200 of the 400 at 100.03, so the total is $10\,002 + 20\,006 = 30\,008$, or 100.0267 a share. The mid-price before was $(99.98 + 100.02)/2$. Afterwards it is $(99.98 + 100.03)/2$.

\iqwrong{Using the best offer for all 300 shares. Only the displayed 100 is available at 100.02.}
\end{solution}

\begin{solution}{iq:iv2:markets-and-microstructure:2}
The mid-price is 100.01. The microprice\index{microprice} is 100.018. You would expect the next move to be up, because little size stands on the offer.

The microprice is $(100.00 \times 100 + 100.02 \times 900)/1000$. The price with the smaller queue is closer to being traded through, so the estimate of fair value leans towards it. Note the crossing of the weights: the large bid size pulls the estimate toward the ask.

\iqwrong{Weighting each price by its own size, which gives 100.002 and leans the wrong way.}
\iqfollow{Is a large bid always good news?}{No. Displayed size can be cancelled, and in a thin market it can be there to mislead. So treat it as a statistic to be tested and not as a fact.}
\end{solution}

\begin{solution}{iq:iv2:markets-and-microstructure:3}
Under price-time priority, $A$ gets 100, $B$ gets 150, and $C$ gets none. Under pro-rata, they get about 71.4, 142.9 and 35.7. Time priority rewards speed, and pro-rata rewards size.

The orders' total is 350. Pro-rata shares are $250 \times 100/350$, $250 \times 200/350$ and $250 \times 50/350$. In practice they round to whole lots, with a rule for the remainder.

Under pro-rata, showing more than you want gets you more of a small trade, so people oversize their quotes. Under time priority, the value lies at the front of the queue, so firms race to get there.

\iqwrong{Thinking that priority rules only matter to brokers. They decide what a market maker's edge is made of.}
\end{solution}

\begin{solution}{iq:iv2:markets-and-microstructure:4}
The effective spread\index{effective spread} is 6 cents, the realised spread\index{realised spread} is 2 cents, and the price impact is 4 cents. The market maker kept the realised spread.

Effective spread is twice the distance from the mid-price at the trade: $2 \times (100.03 - 100.00)$. Realised spread is twice the distance from the later mid-price: $2 \times (100.03 - 100.02)$. Price impact is twice how far the mid-price moved: $2 \times 0.02$. The two parts add up: $2 + 4 = 6$.

You paid 3 cents. The market maker sold at 100.03 what is worth 100.02 five minutes later. So 1 cent of its 3 remains, and the other 2 went to adverse selection.

\iqwrong{Treating the effective spread as profit for the market maker. Most of it can be lost to the price's move.}
\end{solution}

\begin{solution}{iq:iv2:markets-and-microstructure:5}
You save 1.5 cents a share. A market maker nets 1.4 cents on a round trip.

Crossing costs you the 1-cent spread and the 0.3-cent fee, or 1.3 cents more than the bid. Filling passively at the bid earns you the 0.2-cent rebate. The difference is 1.5 cents. A market maker buying at the bid and selling at the offer earns 1 cent and two rebates: $1 + 2 \times 0.2 = 1.4$.

This is a maker-taker\index{maker-taker} venue. The fee is partly in the quoted price, so the spread you see is not what the two sides really pay. Net of fees, the exchange keeps 0.1 cents.

\iqwrong{Assuming you always get the saving. A passive order fills less often, and it fills when the price is about to go through you.}
\end{solution}

\begin{solution}{iq:iv2:markets-and-microstructure:6}
At \$5, one cent is 20 basis points, and at \$500, it is 0.2 basis points. The \$5 stock is more likely to be stuck at one tick with a long queue. There, position in the queue matters more than price, since no one can improve the price.

If the natural spread is below the tick, say 5 basis points, the \$5 stock cannot quote that. So everyone joins one level, and the queue is long. The \$500 stock's spread can vary with little constraint, so its book is thin at each of many levels, and price competition works.

\iqwrong{Thinking low-priced stocks are cheap to trade. In percentage terms, a tick-constrained spread can be much larger than the fair spread.}
\end{solution}

\begin{solution}{iq:iv2:markets-and-microstructure:7}
The spread is 4 cents, from $s = 2\sqrt{-\mathrm{Cov}} = 2\sqrt{0.0004}$. The autocovariance is negative because a buy at the ask tends to be followed by a trade at the bid, which pulls the price back.

This is the Roll\index{Roll model} estimator. The price is $p_t = m_t + \tfrac{s}{2}q_t$, with $q_t = \pm1$ independent. The change is $\Delta m_t + \tfrac{s}{2}(q_t - q_{t-1})$. Only the $q_{t-1}$ term is shared by consecutive changes, so the covariance is $-(s/2)^2\,\mathrm{Var}(q) = -s^2/4$.

With no bounce, $s = 0$, and the covariance is zero, as for a random walk.

\iqwrong{Using it when the autocovariance is positive. The formula then has no real answer, and that tells you the assumptions fail: trades cluster, or the price trends.}
\end{solution}

\begin{solution}{iq:iv2:markets-and-microstructure:8}
The volume-weighted average price\index{VWAP} is 10.11. You beat it by 3 cents, about 30 basis points. The simple average of the three prices is 10.10.

The sum of price times volume is $1\,000 + 2\,040 + 7\,070 = 10\,110$, on 1\,000 shares. Your 10.08 is 3 cents below it, and $0.03/10.11 \approx 0.30\%$.

So the benchmark is weighted by where the volume actually traded. The simple average gives the small print at 10.00 the same weight as the large one at 10.10.

\iqwrong{Thinking that beating the benchmark proves good execution. A buyer who started early may simply have traded before their own impact showed.}
\end{solution}

\begin{solution}{iq:iv2:markets-and-microstructure:9}
Days to cover\index{days to cover} is 15. At 20\% of volume, it would take 75 days. A stock that jumps can force shorts to buy back at any price, and a long exit queue raises their urgency.

The ratio is $30/2 = 15$. Buying 20\% of 2 million is 0.4 million a day, and $30/0.4 = 75$.

The number is the size of a possible forced demand against normal liquidity. Shorts who are stopped out, or whose lender recalls the stock, must buy at any price, and that pushes the price up. It then stops out more of them. This is a short squeeze\index{short squeeze}.

\iqwrong{Saying that a large short interest means the stock must rise. Shorts may be right. It tells you the possible size of the move if they are wrong.}
\end{solution}

\begin{solution}{iq:iv2:markets-and-microstructure:10}
Fair value is \$25.00. The receipt is 80 basis points rich, so sell it, buy the shares and hedge the currency, for about 50 basis points, or about 12.5 cents a receipt.

Two shares cost $\pounds 20$, and at 1.25 that is \$25. Then $0.20/25 = 0.80\%$, and $80 - 30 = 50$ basis points. The conversion between receipt and shares is what ties the two prices, and it costs a fee.

\iqwrong{Forgetting the currency. If you buy the shares in pounds and sell the receipt in dollars, a position remains until you close it or hedge it.}
\end{solution}

\begin{solution}{iq:iv2:markets-and-microstructure:11}
Not under order protection: you should first buy the 100 at 100.01 on $A$, and then the 200 at 100.02 on $B$. A trade-through\index{trade-through} is a trade at a price worse than a protected quote on another venue. An intermarket sweep order lets you send to $B$ at once, on your promise to send to $A$ at the same time.

Rule 611 of Regulation NMS\index{Regulation NMS} makes venues enforce this. An intermarket sweep order is the routing tool: you mark the order to say that you have also taken the better prices, and the venue need not route it elsewhere.

\begin{dated}{2026-11}{Order protection under review}\label{dat:iv2:markets-and-microstructure:rule611}
Rule 611 was adopted in 2005. On 11 June 2026 the SEC proposed to rescind it, together with Rule 610(e). Check where this stands before you state the rule as current.
\end{dated}

\iqwrong{Thinking the best venue is simply where you get the best fill today. The routing obligation sets where you may trade first.}
\end{solution}

\begin{solution}{iq:iv2:markets-and-microstructure:12}
The price is 100, and 300 shares trade. At 100 all the buyers at 100 or better trade, as do all sellers at 100 or lower. The buyers at 99 and the sellers at 101 do not trade.

At each price, compare cumulative demand and supply:
\begin{center}\small
\begin{tabular}{@{}rrrr@{}}
Price & Buyers willing & Sellers willing & Volume \\ \hline
99 & 600 & 150 & 150 \\
100 & 300 & 300 & 300 \\
101 & 100 & 500 & 100
\end{tabular}
\end{center}
At 100 both sides are exactly 300. This is a call auction\index{call auction}: everyone trades at one price.

\iqwrong{Matching orders one at a time from the top, as in a continuous market. The point of the auction is that everybody gets the one price.}
\end{solution}

\begin{solution}{iq:iv2:markets-and-microstructure:13}
You can do about 222\,222 shares, so the order finishes today. That is 11.1\% of the others' volume.

If you are 10\% of the total, then $x = 0.1(2\,000\,000 + x)$, so $x = 200\,000/0.9$. So the 200\,000 can be finished when the others have traded 1.8 million.

A percentage-of-volume order trades more when volume is high, and volume is often high when the price is running. So it tends to buy after rises, and its own trades add to volume, which it then follows.

\iqwrong{Taking 10\% of the others' 2 million, 200\,000. The rate has to include your own trades, or you run at the wrong rate.}
\end{solution}

\begin{solution}{iq:iv2:markets-and-microstructure:14}
Your broker must have reasonable grounds to believe the shares can be borrowed for delivery; this is the locate\index{locate} requirement. A fail to deliver is when the seller does not deliver on settlement date. A short sale without a borrow is not necessarily illegal. Bona fide market makers have an exception from the locate.

Under Regulation SHO\index{Regulation SHO}, Rule 203(b)(1) sets the locate, and Rule 204 sets when a fail must be closed out, by the start of trading on the day after settlement date for a short sale. A market maker must be able to sell before it has found a borrow, in order to quote. The exception does not cover speculative selling.

\iqwrong{Equating a fail with fraud. Fails occur for ordinary operational reasons. The rules work by forcing a close-out.}
\end{solution}

\begin{solution}{iq:iv2:markets-and-microstructure:15}
1\% of a day's volume costs 20 basis points, and 4\% costs 40 basis points. Quadrupling the order multiplies the cost in dollars by 8.

$2\% \times \sqrt{0.01} = 0.2\%$, and $2\% \times \sqrt{0.04} = 0.4\%$. The cost per share doubles, and there are four times as many shares. So the total grows as $Q^{3/2}$.

This is the square-root law\index{square-root impact law} in rough form: an empirical regularity of average impact, with a constant near 1 in many studies. It fits orders well below a whole day's volume, and it is an average. Any given order can do much better or worse.

\iqwrong{Assuming linear impact, so that 4 times the size means 4 times the per-share cost. That overstates large orders, but the total cost is still more than proportional.}
\iqwhy{It sets the capacity of a strategy: the cost per dollar rises with the square root of size, so at some size it eats the edge.}
\end{solution}

\begin{solution}{iq:iv2:markets-and-microstructure:16}
The implementation shortfall\index{implementation shortfall} is \$1\,480, or 29.6 basis points of the \$500\,000 paper portfolio. It is made up of \$800 of execution cost, \$80 of commission and \$600 of opportunity cost.

Execution is $8\,000 \times \$0.10$. Commission is $8\,000 \times \$0.01$. The 2\,000 shares you did not buy would have made $\$0.30$ each. A check by the paper-portfolio route: the paper return is $10\,000 \times 0.30 = 3\,000$. Your real return is $8\,000 \times 0.20 - 80 = 1\,520$. The difference is 1\,480.

\iqwrong{Leaving out the shares you did not buy. A trader who buys nothing has no execution cost, but has missed the whole move.}
\end{solution}

\begin{solution}{iq:iv2:markets-and-microstructure:17}
You are filled on the 21st event. On average that takes 14 seconds, with a standard deviation of about 3.06 seconds. The chance of being filled within 20 seconds is about 96\%.

The wait is the sum of 21 exponential gaps, a gamma distribution, with mean $21/1.5$ and standard deviation $\sqrt{21}/1.5$. Filled by 20 seconds means at least 21 events in 20 seconds, and that count is Poisson with mean 30.

\iqwrong{Using 20 events. You need the 20 ahead removed, and then one more to take you.}
\iqfollow{What does the model leave out?}{Cancels arise anywhere in the queue, and rates move with the market. Most important, your fill arrives when the queue collapses, often as the price goes through you.}
\end{solution}

\begin{solution}{iq:iv2:markets-and-microstructure:18}
At the front you expect to make 0.3 cents a share, and at the back you lose 0.2 cents. So the queue position is worth 0.5 cents a share in this model.

You earn half a spread against the mid-price, 0.5 cents. The fall is one tick, 1 cent, so its expected cost is 0.2 cents at the front and 0.7 at the back. So $0.5 - 0.2 = 0.3$ and $0.5 - 0.7 = -0.2$.

This is why a queue position has value, and why firms invest in speed to get to the front. The same quote at the same price is a good trade at one place in the line and a bad one at another.

\iqwrong{Treating all fills at a price as the same. The difference between them is who you trade with.}
\end{solution}

\begin{solution}{iq:iv2:markets-and-microstructure:19}
The ask is about 108.33 and the bid about 104.09. They are not symmetric around the mean of 106, because the prior is not one half. After a buy and a sell, your probability of 120 is back at 30\%.

A buy is more likely if the value is high: $P(\text{buy} \mid 120) = 0.25 + 0.375 = 0.625$, and $P(\text{buy} \mid 100) = 0.375$. So by Bayes\index{Glosten-Milgrom model}, $P(120 \mid \text{buy}) = 0.1875/(0.1875 + 0.2625) \approx 0.4167$, and the ask is $100 + 20 \times 0.4167 \approx 108.33$. Similarly, $P(120 \mid \text{sell}) = 0.1125/(0.1125 + 0.4375) \approx 0.2045$, and the bid is about 104.09.

The ask is 2.33 above the mean and the bid 1.91 below it. A buy is less likely when the value is more likely to be low, so it carries more news. After a buy and a sell, the likelihood ratio is $(0.625/0.375)(0.375/0.625) = 1$.

\iqwrong{Quoting symmetrically around 106. Each quote must be the expected value given that it gets hit, and the two news amounts differ.}
\iqfollow{What if the informed share rises to 100\%?}{Then a buy means 120 and a sell means 100. Your bid and ask are 100 and 120, and no ordinary trader can trade inside.}
\end{solution}

\begin{solution}{iq:iv2:markets-and-microstructure:20}
The markout\index{markout} at one second is $+0.7$ cents a share, and at one minute it is $-0.3$ cents. On a million shares a day, that is \$7\,000 and $-\$3\,000$. The right horizon is the one at which you can actually flatten your position, and at a minute you are likely to be losing.

Your profit marked at a later mid-price $m$ is $m - 10.00$ plus the rebate. At one second, $0.5 + 0.2 = 0.7$ cents. At one minute, $-0.5 + 0.2 = -0.3$.

The fall from 1 cent to 0.5 to $-0.5$ is adverse selection showing up. So you should check how long you really hold. Then widen on the flow that is costing you, or get out faster, or stop quoting at this point in the queue.

\iqwrong{Reporting the spread captured at the fill, 1.2 cents, as the desk's profit. That is the number that looks best and means least.}
\end{solution}

\begin{solution}{iq:iv2:markets-and-microstructure:21}
It costs you \$75: \$25 of net financing and \$50 of dividend. The stock must fall 75 cents to \$49.25. Your loss has no upper limit.

Proceeds are \$5\,000. Over a quarter of a year, you earn $4\% \times 5\,000/4 = 50$ and pay $6\% \times 5\,000/4 = 75$, a net 25. You must pay the lender the dividend, $100 \times 0.50 = 50$. Each dollar the stock falls earns \$100, so \$75 is 75 cents a share.

As a check, with the fee equal to the interest and no dividend, the cost is zero.

\iqwrong{Forgetting the dividend, or thinking that a short sale earns nothing on its proceeds. The net rate is what matters, and in a hard-to-borrow stock it is strongly negative.}
\end{solution}

\begin{solution}{iq:iv2:markets-and-microstructure:22}
A fall of 10\% or more from yesterday's close, here to \$45 or below, triggers the test. For the rest of today and all of tomorrow, you can sell short only above the national best bid, so at \$44.89 or higher now. If you own the shares, you can still hit the bid.

Rule 201 of Regulation SHO\index{Regulation SHO} is the alternative uptick rule. It does not stop short selling, but it stops short sales from pressing the bid: short orders must rest. So an execution system has to know its state on each stock, and whether each sale is marked long or short.

\iqwrong{Saying short selling is banned while the test holds. Only the price of the sale is limited.}
\end{solution}

\begin{solution}{iq:iv2:markets-and-microstructure:23}
The benchmark is 100.285. Your time-weighted average is 100.425, about 14 basis points above it. A schedule that follows volume would have got the benchmark exactly.

The volume-weighted average is $0.40 \times 100 + 0.15 \times 100.5 + 0.15 \times 101 + 0.30 \times 100.2 = 100.285$. The simple average is 100.425. Its gap of 14 cents is the result of underweighting the cheap first period.

As a limit case, if all four prices were equal, the two averages would be equal and the gap zero. The choice of schedule only matters to the extent that prices and volumes vary together.

\iqwrong{Thinking this says a volume schedule is cheaper. It says only that it tracks the benchmark. On another day, the first period is the dear one.}
\end{solution}

\begin{solution}{iq:iv2:markets-and-microstructure:24}
The day's average is 100.09. Your report says you paid 21 cents above it. The true cost is about 30 cents. If you were all of the volume, the benchmark would equal your price, and your report would say zero.

The average is $0.7 \times 100 + 0.3 \times 100.30 = 100.09$. Your gap is $100.30 - 100.09 = 0.21$, which is your true 0.30 times $(1 - 0.3)$. So the benchmark contains you, and it flatters a large order.

\iqwrong{Using this benchmark for a large order. Compare with the arrival price\index{arrival price} too, since that does not move when you do.}
\iqwhy{An execution desk paid on this benchmark has an incentive to be a bigger part of the volume, which can be the opposite of what the client wants.}
\end{solution}

\begin{solution}{iq:iv2:markets-and-microstructure:25}
You are indifferent at $\lambda \approx 13.3$. Above that, take the dealer's price.

The dealer gives a certain 0.8\%. Working it costs $0.5\% + \lambda \times (1.5\%)^2$ with a penalty. So $0.003 = \lambda \times 0.000225$, which gives $\lambda = 13.33$. For a sense of scale, the 0.3\% you save by working it comes with 1.5\% of risk: 0.2 standard deviations of mean per unit of risk. Over a day or two, that is a strong Sharpe ratio to be getting paid for no view.

Ask also about what the dealer's fill will reveal: is the market aware of the block? Do you hold other names like it? Does your 0.5\% estimate allow for leakage? What is the position's correlation with the rest of your book?

\iqwrong{Comparing only the expected costs. On those alone, working it always wins, and then no one would ever pay for a risk price.}
\end{solution}

\begin{solution}{iq:iv2:markets-and-microstructure:26}
You win with probability 28.125\%, and you expect to make about 0.48 cents a share per attempt.

Take the unit square, in units of 200 microseconds, with your time $Y = 200 + 200U$ and the cancel's $X = 150 + 200W$. You win if $Y < X$, so if $U < W - \tfrac14$. That is a triangle of side $\tfrac34$, and its area is $\tfrac12 \times \tfrac{9}{16} = \tfrac{9}{32} = 0.28125$.

When you win, you make the 2 cents less the 0.3-cent fee, or 1.7 cents. So $0.28125 \times 1.7 \approx 0.478$.

A race like this has no clear winner, since the slower side wins more than a quarter of the time. Reducing the variance of your own latency matters as well as its mean.

\iqwrong{Saying the faster side always wins. Latencies have spread, and a slower median still wins sometimes.}
\end{solution}

\begin{solution}{iq:iv2:markets-and-microstructure:27}
It pays above a fill chance of one third. At 40\%, you gain 0.1 cents a share.

The gain is $p \times 1 - (1 - p) \times 0.5$, which is zero at $p = \tfrac13$. At $p = 0.4$: $0.4 - 0.3 = 0.1$.

What is missing is who you meet in the pool. Fills in a dark pool\index{dark pool} tend to come from those who are happy to sell to you just before a fall, so a mid-price fill is not always a saving. Also missing is the delay, and whether leakage depends on the pool.

\iqwrong{Treating mid-price as a free saving. A trade at mid in the pool can be worse than a trade at the offer, if the price then falls.}
\end{solution}

\begin{solution}{iq:iv2:markets-and-microstructure:28}
It makes 0.5 cents a share. It can pay for the flow because retail orders are, on average, uninformed about short-term moves, so the adverse selection is small.

Half the spread is 1 cent. It gives up 0.2 in improvement, pays 0.1, and loses 0.2 to adverse selection, so $1 - 0.2 - 0.1 - 0.2 = 0.5$.

The exchange's quote has to cover all flow, including the informed. A dealer that can choose which flow it trades with can offer a better price to it. This raises the question of who bears the adverse selection left on the exchange, and whether the broker should send the order to where the client's price is best.

\iqwrong{Saying it ``must be front-running''. The economics work on the flow's low information content alone.}
\end{solution}

\begin{solution}{iq:iv2:markets-and-microstructure:29}
One-minute returns give about 2.08\%, 4\% too high. One-second returns give about 4.77\%, more than twice the truth. You would sample at about five minutes, or use a noise-robust estimator such as two-scale realised volatility.

Each sampled return adds noise variance $s^2/2$, where $s$ is the spread in return terms: $(0.0004)^2/2 = 8 \times 10^{-8}$. Summed over 390 minutes that is 0.0000312, on top of $0.02^2 = 0.0004$. So the variance is 0.000431, and its root is 2.08\%. Over 23\,400 seconds the noise is 0.00187, and the total is 0.00227, a root of 4.77\%.

The true variance does not depend on the sampling rate, and the noise grows in direct proportion to it. So a signature plot of realised volatility against the sampling interval\index{realised volatility} shows it.

\iqwrong{Assuming that more data is always better. More frequent sampling here gives a more precise estimate of the wrong number.}
\end{solution}

\begin{solution}{iq:iv2:markets-and-microstructure:30}
Sell \$15 million of the index. The remaining volatility is about 19.8\%, or about \$1.98 million a year, about \$124\,000 a day.

The beta is $\rho\sigma_s/\sigma_m = 0.75 \times 30/15 = 1.5$. What is left is $\sigma_s\sqrt{1 - \rho^2} = 30\% \times \sqrt{0.4375} \approx 19.84\%$. Per day that is $1.98/16$.

The hedge removes 56\% of the variance, but only a third of the volatility. Note how much stock-specific risk you still run.

\iqwrong{Hedging dollar for dollar, with \$10 million. That leaves a beta of 0.5 to the market and so leaves market risk on the table.}
\end{solution}

\begin{solution}{iq:iv2:markets-and-microstructure:31}
The best $\delta$ is 2 cents, with a fill rate of $e^{-1} \approx 0.37$ a second and a profit of about 0.74 cents a second. Doubling $k$ cuts $\delta$ to 1 cent, and halves the profit rate, while the fill rate stays at 0.37.

Profit rate is $\delta e^{-k\delta}$, and its derivative is zero at $\delta = 1/k$. There, the profit rate is $1/(ke)$. So competition, as higher $k$, narrows the quote and halves the profit, with fills unchanged.

\iqwrong{Quoting as tight as possible to get more fills. Each fill earns less, and the optimum balances the two.}
\end{solution}

\begin{solution}{iq:iv2:markets-and-microstructure:32}
You expect to pay $-p + (1 - p) \times 1.8$ cents from today's mid-price, which is $1.8 - 2.8p$. Posting beats crossing if $p > 0.8/2.8 \approx 28.6\%$. Do not take anything off for the mid-price falling on a fill: you bought at the bid regardless.

Filled, you pay 1 cent below today's mid-price. Not filled, you pay $0.8 + 1 = 1.8$ cents above it. Compare with crossing at $+1$. In general, $p^* = d/(2 + d)$, and at $d = 0$ posting is always better.

You own the stock in every branch, so its later value does not differ between your choices. The adverse selection shows up in the size of $d$: a fill happens more often when the price is falling, and a non-fill costs you more when it is rising. To subtract more would count it twice.

\iqwrong{Counting adverse selection separately on the filled branch. A market maker, who then has to sell, must do this. A buyer who wants to own the stock must not.}
\end{solution}

\begin{solution}{iq:iv2:markets-and-microstructure:33}
VPIN\index{VPIN} is 0.35. It is meant to measure the share of flow that is one-sided, as a proxy for toxic flow.

The imbalances are 20, 40, 0 and 80, with an average of 35 on a bucket of 100.

The cautions are these. Trade classification is noisy. Imbalance may reflect a large uninformed order, such as an index fund. High readings are also found alongside volatility, so check whether it predicts anything beyond volatility and volume before you size a decision on it.

\iqwrong{Reading a high value as proof of informed trading. It measures imbalance, and the rest is interpretation.}
\end{solution}

\begin{solution}{iq:iv2:markets-and-microstructure:34}
A single sweep costs \$55. The slices cost \$19, and the risk is a standard deviation of about \$84. So the slices save \$36 in expected cost and take on about 2.3 times that in risk.

A sweep pays $1, 2, \ldots, 10$ cents on each 100 shares: $100 \times 55$ cents. The slices pay $1 + 0.2(k - 1)$ cents on slice $k$, for an average of 1.9 cents on 1\,000 shares.

The risk comes from the shares you have yet to buy: 900 after the first minute, then 800, and so on down to 100. The squares add to $100^2 \times 285$, so the root is $100 \times 16.9$, and at 5 cents a minute that is about \$84.

\iqwrong{Ignoring the risk because it has zero mean. A mean-variance view weighs it, and if you have a short-term view, the cost is real.}
\end{solution}

\begin{solution}{iq:iv2:markets-and-microstructure:35}
Sell 625 first and 375 second, at a total cost of about \$143.75. An equal split costs \$150, and all at once costs \$200. At $c = 0$ the best split is equal. At $c = 0$ and $b = 2a$, every split costs the same.

Let $y = x_2$ and $x_1 = X - y$. The cost is $a(X - y)^2 + ay^2 + b(X - y)y + cy^2$. Its derivative in $y$ is $-2a(X - y) + 2ay + b(X - 2y) + 2cy$, which vanishes at
\[
y = \frac{X(2a - b)}{4a - 2b + 2c} = \frac{1000 \times 0.0003}{0.0008} = 375.
\]
Then the cost is $0.0002(625^2 + 375^2) + 0.0001 \times 625 \times 375 + 0.0001 \times 375^2 = 106.25 + 23.44 + 14.06$.

With $c = 0$, $y = X/2$: splitting evenly spreads the convex impact, and risk pushes the split forward. With $c = 0$ and $b = 2a$, the cost is $a(x_1 + x_2)^2$ for every split, so the timing is worthless, and the price effect of the first slice cancels the gain from splitting.

\iqwrong{Thinking risk aversion makes you trade later. Holding is the risk, so it brings selling forward.}
\end{solution}

\begin{solution}{iq:iv2:markets-and-microstructure:36}
The chance the offer goes first is about 91\%, so the expected move is about $+0.41$ of a tick. The microprice puts fair value 0.2 of a tick above the mid-price. The two differ because the microprice is a simple linear weighting, while the race is a tail probability.

The offer goes first if 3 buys arrive before 7 sells. Look at the first 9 orders, which must decide it: the offer goes first if at least 3 are buys. Then $P = 1 - (1 + 9 + 36)/512 = 466/512 \approx 91.0\%$. The expected move is $\tfrac12(0.910 - 0.090) = 0.41$ ticks.

The microprice\index{microprice} weights are $7/10$ and $3/10$, so it lies $0.7 - 0.5 = 0.2$ of a tick above the mid-price. As a check, with equal queues, both say 0.

The race overstates what you would see in practice, because queues are refilled. The truth is usually between the two. So fit the move as a function of imbalance from data, and you usually find that it curves.

\iqwrong{Giving the 70\% from the sizes as the probability. That is the size share, and not the chance of winning a race to depletion.}
\end{solution}

\begin{solution}{iq:iv2:markets-and-microstructure:37}
A buy following a buy moves the price 0.2 cents, and a sell following a buy moves it $-1.8$ cents. After a run of five buys starting from neutral, the price has moved 1.8 cents. Order flow is predictable, and so the price already includes what is expected: only the surprise moves it.

The expected next sign after a buy is 0.8. A buy is a surprise of 0.2, and a sell a surprise of $-1.8$. Each such move has a mean of zero, since $0.9 \times 0.2 + 0.1 \times (-1.8) = 0$, with the buy probability $(1 + 0.8)/2 = 0.9$. So price changes have no autocorrelation, though the signs do.

Over a run of five, the first is a surprise of 1 and the other four 0.2 each, so $1 + 0.8$. The impact of a sequence of buys is concave in length, and that is the same concavity seen in large orders. This is a simple form of the Madhavan-Richardson-Roomans model\index{Madhavan-Richardson-Roomans model}.

\iqwrong{Saying that buy pressure must keep moving the price up because orders are autocorrelated. Prices that could be predicted so easily would be traded away.}
\iqfollow{And if you are the one splitting the order?}{The market is reading your run as a probable continuation, and so it prices you in. That makes your first slices dear and your later ones cheap.}
\end{solution}

\begin{solution}{iq:iv2:markets-and-microstructure:38}
It takes 2.5 days. Your impact cost is about 0.94\%. The standard deviation is about 1.83\%. So you quote a discount of about 1.86\%. Hedging with an index cuts the standard deviation to about 1.10\%, and half of that brings the discount to about 1.49\%.

The rate is $0.2 \times 2$ million a day, so 1 million takes 2.5 days. Impact is $\tfrac23 \times 2\% \times \sqrt{0.5} \approx 0.943\%$.

For the risk, you hold a share falling linearly from 1 to 0 over $T$ days. Your variance is $\sigma^2\int_0^T (1 - t/T)^2 dt = \sigma^2T/3$. So the standard deviation is $2\% \times \sqrt{2.5/3} \approx 1.83\%$. With a hedge, $1.2\% \times \sqrt{2.5/3} \approx 1.10\%$.

The total is $0.943 + \tfrac12 \times 1.826 \approx 1.86\%$. With the hedge it is $0.943 + \tfrac12 \times 1.095 \approx 1.49\%$. Then you add the risk that this client knows something: if the client is informed, you lose whatever the news is.

\iqwrong{Charging only the impact. The dealer has no view on the stock, and for the 2.5 days of risk, it must be paid.}
\end{solution}

\begin{solution}{iq:iv2:markets-and-microstructure:39}
The bid-side indifference price is $s - \gamma\sigma^2\tau(q + \tfrac12)$, and the ask-side is $s - \gamma\sigma^2\tau(q - \tfrac12)$. Their average, the reservation price\index{reservation price}, is $s - q\gamma\sigma^2\tau$. So with these numbers it is 99, and you quote about 98.25 and 99.75.

With exponential utility and a normal price, the certainty equivalent of holding $q$ shares is $qs - \tfrac12\gamma q^2\sigma^2\tau$. Buying one more at $p$ is worth it if $(q + 1)s - \tfrac12\gamma(q + 1)^2\sigma^2\tau - p \geq qs - \tfrac12\gamma q^2\sigma^2\tau$. That gives $p \leq s - \gamma\sigma^2\tau(q + \tfrac12)$. Selling one gives the other in the same way.

With the numbers, $\gamma\sigma^2\tau = 0.1 \times 4 \times 0.5 = 0.2$, and the reservation price is $100 - 5 \times 0.2 = 99$. The spread is $0.2 + 20\ln(1 + 0.1/1.5) \approx 1.49$, so half of it is 0.745. These are the Avellaneda-Stoikov\index{Avellaneda-Stoikov model} quotes.

If $q = 0$, the quotes centre on $s$. If you are long, you lean them down to sell and to buy less.

\iqwrong{Skewing only one side. Both quotes shift by the same amount, and that keeps the same spread with a new centre.}
\end{solution}

\begin{solution}{iq:iv2:markets-and-microstructure:40}
No. The first is spoofing\index{spoofing}: bidding with the intent to cancel before execution. The second can be a lawful strategy whose orders do change. US law makes the first illegal through its intent. The Commodity Exchange Act, at 7 U.S.C. 6c(a)(5)(C), prohibits placing a bid or offer with the intent of cancelling it before it executes.

So a cancel rate cannot be what decides it. In the case against one trader, a US appeals court rejected the argument that the rule is vague because high-frequency traders cancel 98\% of their orders. The issue is what the order was for. The first algorithm's large bids are placed only to move others' prices, so they are placed not to trade, while the second's were placed to trade at the price then shown.

To detect it, look at orders where large resting size on one side coincides with fills on the other, and where the large orders are cancelled as soon as those fills occur. Compare the distance of the orders from the market, the order-to-trade ratio per pattern, and the share of orders on the far side that ever fill. Escalate to compliance when you find a pattern. Do not judge intent by yourself.

\iqwrong{``We cancel everything, so we are all spoofers'' or ``we did not intend harm, so it is fine''. The law looks for an intent to cancel at the time of the order, and it does not need a harm.}
\iqfollow{Is this only for futures?}{The cited provision is for markets under that Act, and securities and other markets have their own manipulation rules. Know which regime covers each book.}
\end{solution}

\begin{solution}{iq:iv2:markets-and-microstructure:41}
My recommended answer: long-term investors with small orders probably gain from a narrower spread, those with large orders may not, front-of-queue market makers lose, back-of-queue ones may gain, and fast traders who can step ahead gain. The net effect depends on how much the natural spread is below the tick.

The tick is a floor on the spread. If the natural spread is below 1 cent, halving the tick narrows it. A small order then pays less.

But a tick also protects a queue. Today, to get ahead of a long queue, you must improve by a whole cent, and that is expensive. With a half-cent tick, stepping ahead costs half as much. So the large displayed size that sits there today becomes less worthwhile, and depth falls at the best quote. A large order may now have to pay for more levels, so it may gain less, or lose.

Those at the front of the queue lose a position they paid for. Those at the back find the queue shorter. Fast traders can reprice for cheap, and they gain the most in stepping ahead.

To judge it, measure quoted and effective spreads, depth at and near the best quote, the cost of large orders through implementation shortfall, realised spreads and queue lengths, and the market maker's adverse selection. Compare with similar stocks not changed, before and after.

\iqwrong{``A smaller tick means a smaller spread, so everybody gains.'' Spread is one cost, and depth and queue value change too.}
\end{solution}

\begin{solution}{iq:iv2:markets-and-microstructure:42}
The break-even half-spread is 0.5 cents. With one firm, it is 0. With very many, it tends to about 0.67 cents. Removing the race would remove this component of the spread.

Ordinary trades pay $h$ at a rate of 1 a second. A signal comes at a rate of 0.5. With three firms, the quoter is first to cancel one time in three, and so it is sniped two times in three. Each time it loses $2 - h$. So
\[
1 \times h = 0.5 \times \tfrac23 \times (2 - h),
\]
which gives $h = \tfrac23/\tfrac43 = 0.5$. With $N$ firms, the factor is $(N - 1)/N$: at $N = 1$ the term vanishes, and as $N$ grows it tends to 1, so $h \to 1/1.5$.

This is the result in Budish, Cramton and Shim's work on the arms race\index{latency arbitrage}: the race itself is the cost, and it is paid for in the spread by the ordinary traders. Frequent batch auctions were proposed to remove it. Being faster does not reduce it: if everyone is faster, the three-way tie remains.

\iqwrong{Saying that latency investment narrows spreads because prices update faster. Here $h$ does not depend on how fast the firms are, only on how many there are.}
\end{solution}

\begin{solution}{iq:iv2:markets-and-microstructure:43}
You expect to make $50q/(1000 + q)$ per trade. It rises with $q$ and flattens towards 50. If you may take at most 200 lots in one trade, show about 333 lots.

Your share of any trade is $q/(1000 + q)$. On a small trade, you earn half a tick on 200 lots. On a large one, you earn half a tick and lose a tick on 800 lots. So the amount per unit share is $0.9 \times 200 \times 0.5 + 0.1 \times 800 \times (0.5 - 1) = 90 - 40 = 50$.

So at $q = 100$ you make about 4.5, at $q = 1\,000$ you make 25, and at $q = 3\,000$ you make 37.5. More size always helps, in expectation, so everyone oversizes. What stops you is the risk: your largest fill is $800q/(1000 + q)$, and setting that equal to 200 gives $q = 1000/3$.

And others see the same thing, so their size grows too, and the 1\,000 is not fixed. This is why pro-rata books show far more than people want.

\iqwrong{Showing what you want to trade. In pro-rata you are filled for only part of your display, so you show more, and you risk being filled for all of it.}
\end{solution}

\begin{solution}{iq:iv2:markets-and-microstructure:44}
It is not significant: the standard error of the difference is about 2.12 basis points, so $t \approx 0.57$. You would need about 9\,800 orders per broker. The cost is so noisy relative to the gap that you cannot rank brokers on a few hundred orders.

The standard error is $30 \times \sqrt{2/400}$. For the required sample, $n = 2(1.96 + 0.84)^2(30/1.2)^2 = 2 \times 7.84 \times 625 = 9\,800$.

Two ways to do better. Use a control: comparing each order's cost with a model of what it should cost, given size, volatility and spread, can take out much of the variance. And balance the orders by difficulty, so that luck in the mix does not decide it. The random assignment in the stem is the valuable part, because without it, one broker may get the hard orders.

\iqwrong{Choosing the cheaper broker on 400 orders, and then being surprised next quarter. Over a sample like this, the rank is close to a coin toss.}
\end{solution}

\begin{solution}{iq:iv2:markets-and-microstructure:45}
The predator's best $y$ is 25\,000 shares, half a slice, for a profit of \$625. Your extra cost is \$2\,500, which is a third of your cost of \$7\,500 with no predator. So a predictable schedule costs you much more than it makes for the predator, and you should randomise your timing and size and hide your total.

With no predator, you sell $x = 50\,000$ at $p_0 - bx$, and then $x$ at $p_0 - 2bx$. So your cost is $3bx^2 = 3 \times 2\,500 = 7\,500$, since $bx^2 = 2\,500$.

The predator sells $y$ at $p_0 - b(x + y)$, and buys back at $p_0 - 2bx$. The profit is $y \cdot b(x - y)$, and it is largest at $y = x/2$, where it is $bx^2/4 = 625$.

You receive less by $by$ on each of your slices, since the predator's $y$ stays in the price until the buy-back. That is $2bxy = bx^2 = 2\,500$. The rest of your loss goes to those who bought from you and from the predator at the lower prices.

\iqwrong{Thinking that a time-weighted order has no leakage because each slice is small. A predictable pattern is itself the leak.}
\end{solution}

\begin{solution}{iq:iv2:markets-and-microstructure:46}
The best horizon is about 0.58 days. There, cost and risk are each about 0.88\%, for a total of about 1.75\%. The empirical square-root law says that cost depends on $Q/V$ with little dependence on time, so this model's gain from going slowly may be overstated.

Write the objective as $aT^{-1/2} + bT^{1/2}$, with $a = \tfrac23 \times 2\% \times \sqrt{0.25} \approx 0.667\%$ and $b = 2\%/\sqrt 3 \approx 1.155\%$. The derivative gives $T = a/b \approx 0.577$. At that point the two terms are equal, each $\sqrt{ab} \approx 0.877\%$.

Where the law is true, cost barely falls with a longer horizon, as long as the orders stay in its range. Then going slowly is only risk, and the right answer is shorter. Check which applies in your own data: bucket the orders by duration, and see whether the cost at a given $Q/V$ depends on it.

\iqwrong{Saying that slower is always cheaper. Slower adds price risk and leakage. And the data do not show large gains from stretching a fixed order.}
\end{solution}

\begin{solution}{iq:iv2:markets-and-microstructure:47}
The bands are at \$95 and \$105. If the market stays against the band for 15 seconds in a limit state, the listing exchange declares a pause of at least 5 minutes. The bands are doubled in the first 15 minutes of the day, and in the 25 minutes before the close. During the pause you cannot hedge your delta in this stock, and you are exposed. So you hedge with something related that still trades, or size down.

Under the limit up-limit down\index{limit up-limit down} plan, the bands are set around the mean trade price of the past five minutes. A trade cannot print outside the band. In a limit state, the offer is at the lower band and there is no bid there, so sellers cannot find buyers at an allowed price, and the price cannot fall further for now. At the end of 15 seconds, it pauses. The reopening uses an auction.

You sold a put, and hedged by being short stock. As the stock falls, the put's delta grows and you must short more. In the pause you cannot, and the reopening can gap well below the band, so you take a gamma loss you could not hedge against.

Use index futures or a sector fund as a partial hedge, and accept the basis. Check how your option venue treats the pause before you rely on it. A reasonable desk also limits short gamma in names most likely to halt, such as those with earnings.

\iqwrong{Thinking the bands protect you from a gap. They stop trading at a price. They do not stop the price from opening lower.}
\end{solution}

\begin{solution}{iq:iv2:markets-and-microstructure:48}
The insider trades $x = \beta(v - p_0)$ with $\beta = \sigma_u/\sigma_v$, and the price responds with $\lambda = \sigma_v/(2\sigma_u)$. With the numbers, $\beta = 500$ shares per dollar, $\lambda = 0.001$ dollars per share, the insider's expected profit is \$1\,000, and the remaining variance is 2. The noise traders pay.

Given $v$, the insider chooses $x$ to maximise $E[x(v - p_0 - \lambda(x + u))] = x(v - p_0) - \lambda x^2$. So $x = (v - p_0)/(2\lambda)$ and $\beta = 1/(2\lambda)$.

The market maker prices at the regression of $v$ on the flow $y = x + u$: $\lambda = \mathrm{Cov}(v, y)/\mathrm{Var}(y) = \beta\sigma_v^2/(\beta^2\sigma_v^2 + \sigma_u^2)$. Substituting $\beta = 1/(2\lambda)$ gives $\beta^2\sigma_v^2 = \sigma_u^2$. So the insider's order has the same variance as the noise, and $\lambda = \sigma_v/(2\sigma_u)$.

The expected profit is $E[(v - p_0)^2]/(4\lambda) = \sigma_v\sigma_u/2 = 2 \times 1000/2 = 1\,000$. The variance of $v$ given $y$ is $\sigma_v^2(1 - \tfrac12) = 2$: half the information enters the price.

In this Kyle model\index{Kyle model} the market maker breaks even, on average. So the insider's profit equals the noise traders' loss. Doubling the noise halves $\lambda$ and doubles the insider's profit, while the price remains as informative.

\iqwrong{Thinking more noise makes the market less informative. The insider simply trades more, so that exactly half the variance always goes into the price.}
\iqwhy{This is where $\lambda$, the price change per share, comes from as a measure of liquidity, and why you read it as a statement about information.}
\end{solution}

\begin{solution}{iq:iv2:markets-and-microstructure:49}
The optimal path is $x(t) = X\sinh(\kappa(T - t))/\sinh(\kappa T)$. By half way you have sold about 55.7\% at $\kappa T = 1$, and about 78.7\% at $\kappa T = 3$. Against selling at a constant rate, the cost rises about 1.9\% and the variance falls about 11.7\% at $\kappa T = 1$. At $\kappa T = 3$, the cost rises about 55\%, and the variance falls about 51\%. At $\lambda = 0$, the path is a straight line.

With $v = -\dot x$, the objective's integrand is $L = \eta\dot x^2 + \lambda\sigma^2x^2$. The Euler-Lagrange\index{calculus of variations} equation $2\eta\ddot x = 2\lambda\sigma^2 x$ gives $\ddot x = \kappa^2x$. With $x(0) = X$ and $x(T) = 0$, you get the $\sinh$ form above. As $\kappa \to 0$, it tends to $X(1 - t/T)$.

The share sold by half way is $1 - \sinh(\kappa T/2)/\sinh(\kappa T)$: $1 - 0.5211/1.1752$ and $1 - 2.1293/10.018$.

With $T = 1$, integrating gives
\[
\int_0^1\dot x^2\,dt = X^2\kappa^2\frac{\tfrac12 + \sinh(2\kappa)/(4\kappa)}{\sinh^2\kappa},\qquad
\int_0^1 x^2\,dt = X^2\frac{\sinh(2\kappa)/(4\kappa) - \tfrac12}{\sinh^2\kappa}.
\]
For constant selling, these are $X^2$ and $X^2/3$. So at $\kappa = 1$: 1.019 and 0.2945. At $\kappa = 3$: 1.552 and 0.1625.

This is the Almgren-Chriss\index{Almgren-Chriss model} model. The path is steep at the start because holding is the risk. The trade-off is shown best on the efficient frontier: a small move from linear gets a lot of risk reduction for little cost, and later moves are expensive.

\iqwrong{Selling evenly and calling that the optimum. It is optimal only with no risk aversion. Any $\lambda > 0$ brings selling forward.}
\iqfollow{What does $\kappa$ tell you?}{Its inverse is the natural time of the order: about how long an order takes for a trader with this cost and risk. If $T$ is much above it, most of the order goes out early.}
\end{solution}

\begin{solution}{iq:iv2:markets-and-microstructure:50}
The sizes are in proportion to $1$, $1 - a$ and $1$. With $a = 0.5$, the shares are 40\%, 20\% and 40\%, at a cost of $0.3\kappa X^2$. That is about 1.8\% below equal thirds, at about $0.3056\kappa X^2$, and 40\% below all at once, at $0.5\kappa X^2$. As $a \to 0$, the sizes tend to equal thirds. As $a \to 1$, the middle trade vanishes, and you trade half at each end.

Write the cost as $\tfrac12\kappa\,x^\top Gx$, with $G$ having 1 on the diagonal, $a$ next to it and $a^2$ at the corners. The optimum under $\sum_i x_i = X$ has $Gx$ constant. Try $w = (1, 1 - a, 1)$. The first row gives $1 + a(1 - a) + a^2 = 1 + a$. The middle row gives $a + (1 - a) + a = 1 + a$. So $Gw = (1 + a)\mathbf 1$, and the optimum is proportional to $w$.

With the total $S = 3 - a$, the cost is $\tfrac12\kappa X^2(1 + a)/S$. At $a = 0.5$: $\tfrac12 \times 1.5/2.5 = 0.3$. Equal thirds cost $\tfrac12\kappa(X/3)^2(3 + 4a + 2a^2) = \tfrac12\kappa X^2 \times 5.5/9$.

This is the shape in the Obizhaeva-Wang\index{Obizhaeva-Wang model} model: large trades at the start and end, and a slower trickle between. The trade at the start has the whole time to decay. The trade at the end leaves nothing behind to hurt, so it can be large.

\iqwrong{Assuming a smooth schedule must be best. With a book that recovers, the first and last trades are special, and the model says to take advantage of it.}
\iqfollow{What if the decay is not exponential?}{The result depends on the shape of decay. With a power-law decay, or if impact is non-linear, the schedule differs, and some shapes allow a profitable round trip, a sign that the model is not sound.}
\end{solution}
